\documentclass[a4paper,10pt,twoside,openany]{article}

\usepackage[lang=hebrew]{maths}
\usepackage{hebrewdoc}
\usepackage{stylish}
\usepackage{lipsum}
\let\bs\blacksquare

\usepackage{siunitx,array}
\sisetup{table-format=-1.5, group-digits=false}
\newcolumntype{L}{>{\phantom{$-$}}l}       % prefix some whitespace
\newcommand\mcL[1]{\multicolumn{1}{L}{#1}} % handy shortcut macro

\setlength{\parindent}{0pt}

%%%%%%%%%%%%
% Styling %
%%%%%%%%%%%%

\usepackage{enumitem}

%%%%%%%%%%%%%
% Counters  %
%%%%%%%%%%%%%

\setcounter{section}{1}     
            
%BIBLIOGRAPHY
\usepackage[
backend=biber,
style=alphabetic,
]{biblatex}
\addbibresource{bibliography.bib} %Imports bibliography file

%%%%%%%%%%
% Title  %
%%%%%%%%%%
\title{
אלגברה ב' - גיליון תרגילי בית 5 \\
משפט קיילי-המילטון, סכומים ישרים, והטלות
\\
\vspace{1cm}
\large{תאריך הגשה: 4.6.2026}
}
\date{}

\begin{document}
\maketitle

\begin{exercise}%1
יהי
$V$
מרחב וקטורי ממימד
$n \in \mbb{N}_+$
מעל שדה
$\mbb{F}$
ותהי
$T \in \endo_{\mbb{F}}\prs{V}$.

\begin{enumerate}
\item
יהי
\[\text{.} \mbb{F}\brs{T} \ceq \set{p\prs{T}}{p \in \mbb{F}\brs{x}} \leq \endo_{\mbb{F}}\prs{V}\]
 הראו כי
\[\text{.} \mbb{F}\brs{T} = \spn\prs{I, T, \ldots, T^{n-1}}\]

\item
מצאו דוגמא עבור
$S \in \endo_{\mbb{F}}\prs{V}$
עבורה
\[\set{I, S, \ldots, S^{n-1}}\]
קבוצה בלתי־תלויה לינארית. הסיקו שבמקרה זה
\[\text{.} \dim_{\mbb{F}} \mbb{F}\brs{S} = n\]

\item
נניח כי
$T$
הפיכה. הראו שמתקיים
$T^{-1} \in \mbb{F}\brs{T}$.

\item נניח כי
$T$
הפיכה ונסמן
\[\text{.} \mbb{F}\prs{T} \ceq \set{\sum_{i=-m}^m a_i T^i}{\substack{m \in \mbb{N} \\ \forall i: a_i \in \mbb{F}}} \leq \endo_{\mbb{F}}\prs{V}\]
הראו כי לכל
$S_1, S_2 \in \mbb{F}\brs{T}$
מתקיים
$S_1 \circ S_2 \in \mbb{F}\brs{T}$
והסיקו שמתקיים
$\mbb{F}\brs{T} = \mbb{F}\prs{T}$.
\end{enumerate}
\end{exercise}

\begin{exercise} %2
\begin{enumerate}
\item תהי
$T \in \endo_{\mbb{C}}\prs{\mbb{C}^4}$
שהערכים העצמיים שלה הם
$3,5,8$.
הראו כי
\[\text{.} \prs{T - 3I}^2\prs{T - 5I}^2 \prs{T - 8I}^2 = 0\]
\item תהי
$T \in \endo_{\mbb{C}}\prs{V}$
הטלה ויהיו
\[\text{.} m \ceq \dim \ker \prs{T}, \quad n \ceq \dim \im \prs{T}\]
הראו כי
\[\text{.} p_T\prs{x} = x^m \prs{x-1}^n\]
\end{enumerate}
\end{exercise}

\begin{exercise} %3
יהי
$V = \Mat_n\prs{\RR}$
ויהיו
\begin{align*}
V_{\mrm{sym}} &\coloneqq \set{A \in V}{A^t = A} \\
\text{.} V_{\mrm{antisym}} &\coloneqq \set{p \in V}{A^t = -A}
\end{align*}
\begin{enumerate}
\item
הראו כי
\[\text{.}V = V_{\mrm{sym}} \oplus V_{\mrm{antisym}}\]
\item
מיצאו את ההטלה על
$V_{\mrm{sym}}$
במקביל ל־%
$V_{\mrm{antisym}}$.
\item
נניח כי
$n = 2$,
ויהי
\[\mcal{E} \coloneqq \prs{E_{1,1}, E_{1,2}, E_{2,1}, E_{2,2}}\]
הבסיס הסטנדרטי של
$V$.
מיצאו את המטריצה המייצגת של ההטלה מהסעיף הקודם לפי
$\mcal{E}$.
\end{enumerate}
\end{exercise}

\end{document}